\begin{trapbox}
	\begin{description}[style=nextline, leftmargin=2.8em, labelsep=0.5em, itemsep=0.35em]
		\item[\textbf{\textcolor{CTrap}{易错点一（方向混淆）}}] 误将比例记成 \\(AG:GD=1:2\\)，即弄反了顶点侧和底边侧。
		\item[\textbf{\textcolor{CTrap}{正确理解（方向纠正）：}}] 靠近顶点的一段较长，正确比例为 \\(AG:GD=2:1\\)。
		\item[\textbf{\textcolor{CTrap}{错因分析（记忆偏差）：}}] 对结论字面记忆不牢，或受其他比例（如黄金分割）干扰。
	\end{description}

	\medskip
	\begin{description}[style=nextline, leftmargin=2.8em, labelsep=0.5em, itemsep=0.35em]
		\item[\textbf{\textcolor{CTrap}{易错点二（特殊化错误）}}] 认为该比例只在等边三角形或特定三角形中成立。
		\item[\textbf{\textcolor{CTrap}{正确理解（普遍成立）：}}] 重心分中线 \\(2:1\\) 对任意三角形都成立，与形状无关。
		\item[\textbf{\textcolor{CTrap}{错因分析（归纳不完整）：}}] 只记忆了等边三角形的特例，忽略了重心性质的普适性。
	\end{description}

	\medskip
	\begin{description}[style=nextline, leftmargin=2.8em, labelsep=0.5em, itemsep=0.35em]
		\item[\textbf{\textcolor{CTrap}{易错点三（整体关系误用）}}] 直接认为 \\(AG = \frac12 AD\\) 或 \\(GD = \frac12 AD\\)。
		\item[\textbf{\textcolor{CTrap}{正确理解（分数关系）：}}] \\(AG = \frac23 AD\\)，\\(GD = \frac13 AD\\)，即整体被分为三份。
		\item[\textbf{\textcolor{CTrap}{错因分析（比例与分数混淆）：}}] 从比例 \\(2:1\\) 直接认为 \\(AG\\) 占一半，其实 \\(AG\\) 占整体的 \\(\frac23\\)。
	\end{description}
\end{trapbox}
